\clearpage
\fancyhead[L]{\small JUNO.}
\hypertarget{native-261-lower}{}
\begin{center}
JUNO.

\rule{0.75\linewidth}{0.4pt}

Monatliche Correspondenz. Band X. S. 463–471. 1804 November.
\end{center}

\noindent—\quad—\quad—\par
Dr. \textsc{Gauss} verlangte von uns die Bestimmung eines Sterns, welcher am 28. Sept. sehr nahe bei der \emph{Juno} stand, und mit welchem er diese verglichen hatte. Wir fanden aus einem Mittel fünftägiger Beobachtungen folgende Position dieses Sterns, welche auch andern Beobachtern dienen kann.

\begin{center}
\begingroup\fontsize{9}{11}\selectfont\setlength{\tabcolsep}{3pt}\renewcommand{\arraystretch}{1.1}
\begin{tabular}{r|r|r|r|r}
Grösse & \shortstack{Mittlere gerade\\Aufsteigung\\1804} & \shortstack{Jährliche\\Veränderung} & \shortstack{Mittl. südliche\\Abweichung\\1804} & \shortstack{Jährliche\\Veränderung}\\ \hline
8 & $358^\circ\ 19'\ 34''42$ & $+46''1$ & $5^\circ\ 13'\ 57''27$ & $-20''0$ \\
\end{tabular}\endgroup
\end{center}


Unermüdet fährt Dr. \textsc{Gauss} in seiner beschwerlichen und mühevollen Arbeit fort, die Elemente dieser neuen Planeten-Bahn zu verbessern. „Seit meinem letzten Briefe“, (schreibt dieser grosse Calculator) „habe ich zwar unsere \emph{Juno} noch immer fleissig beobachtet, ich mache aber von meinen eigenen Beobachtungen für die Elemente keinen Gebrauch mehr, da die Ihrigen, mir von ihrer Güte mitgetheilten nun schon einen ansehnlichen Bogen zu befassen anfangen. Inzwischen habe ich nach einem vergeblichen Versuche, bei einer neuen Verbesserung der Juno-Bahn die mittlere Bewegung der Ceres und Pallas auch für diese hypothetisch zum Grunde zu legen, eine zweite, von Hypothesen unabhängige Bestimmung der Elemente gemacht, wovon dieses die Resultate sind :“

\begin{center}
\emph{II\textsuperscript{te} Elemente der Juno.}
\end{center}

\begin{center}
\begingroup\fontsize{9}{11}\selectfont\setlength{\tabcolsep}{3pt}\renewcommand{\arraystretch}{1.1}
\begin{tabular}{lr}
Epoche 1804 Sept. 30. Um $0^{\mathrm u}$ mittl. Zeit  &  \\
\quad in Seeberg . . . . . . . . . & $21^\circ\ 17'\ 47''$ \\
Tägliche mittlere Bewegung . . . . & $836''89$ \\
Sonnenferne . . . . . . . . . & $231^\circ\ 38'\ 1''$ \\
Excentricität . . . . . . . . . & $0.254964$ \\
Logar. der halben Achse . . . . . & $0.4182255$ \\
Knoten . . . . . . . . . . . & $170^\circ\ 46'\ 41''$ \\
Neigung der Bahn . . . . . . . & $12^\circ\ 19'\ 43''$ \\
\end{tabular}\endgroup
\end{center}


„Sind in meine Rechnungen keine Fehler eingeschlichen, so darf man es schon als ziemlich ausgemacht ansehen, dass die mittlere Bewegung der Juno der der Ceres
 und Pallas \emph{nicht gleich}, sondern \emph{beträchtlich grösser} ist, mithin Umlaufszeit und mittlere Entfernung von der Sonne kleiner. Die fortgesetzten Beobachtungen werden uns sehr bald Gewissheit darüber geben.“

Anfangs schien es, als ob Dr. \textsc{Gauss} auch für die \emph{Juno} dieselbe Umlaufszeit wie für die Ceres und Pallas finden würde. Dies würde Dr. \textsc{Olbers} Hypothese über die Entstehung dieses kleinen Plane%
\SourceLeaf{262}%
ten sehr zuwider gewesen sein; denn, ob sich gleich zeigen lässt, dass, wenn man voraussetzt, alle diese
kleinen Planeten könnten vielleicht nur Trümmer und Bruchstücke eines zerstörten grössern Planeten
sein, die Umlaufszeiten dieser kleinen Trümmer nicht \emph{sehr ungleich} sein können, so war es doch sehr
unwahrscheinlich anzunehmen, dass alle die Stücke des zersprengten Planeten genau dieselbe Geschwindigkeit erhalten haben sollten; aber nun findet unser vortrefflicher \textsc{Gauss}, wie wir oben gesehen haben,
die mittlere Bewegung der \emph{Juno} viel schneller als die der Ceres und Pallas.

Dr. \textsc{Olbers} drückt sich hierüber folgendermaassen aus: „Die ganze Lage der Juno-Bahn hat nichts,
was nicht mit meiner Hypothese (die ich übrigens auch noch für weiter nichts als eine Hypothese ausgeben will) zu vereinigen wäre; ihre Knoten mit der Ceres-Bahn fallen jetzt etwa 24 Grade von dem
Knoten der Pallas-Bahn, allein bei den schon so verschiedenen Neigungen dieser Bahnen müssen sich
die Knoten durch die anziehende Kraft des Jupiter ungleichförmig verrücken. Jetzt liegt die Juno-Bahn beim niedersteigenden Knoten auf der Ceres-Bahn, bei der die Pallas-Bahn dieser so nahe ist,
weit innerhalb der Ceres-Bahn; aber da die Aphelien aller dieser Bahnen eine ganz andere Bewegung
haben, als die Knoten, die Lagen der Apsiden-Linien gegen die Knoten sich also immer verändern, und
da diese Bahnen fast gleich grosse Achsen, aber sehr ungleiche Excentricitäten haben, so folgt, dass
sich diese Bahnen zu gewissen Zeiten wirklich schneiden werden, und auch in ehemaligen Zeiten wirklich geschnitten haben. Nehme ich z.\,B. die von \textsc{Oriani} bestimmte jährliche Verrückung der Aphelien
für die Pallas $106''1$ und für die Ceres $120''9$ an, und setze die Knoten als siderisch ruhend und die
Neigungen unveränderlich, so folgt, dass sich die Bahnen der Ceres und Pallas beim niedersteigenden
Knoten der Pallas auf der Ceres-Bahn vor 7463 Jahren wirklich geschnitten und nach 282 Jahren wieder
schneiden werden; beim aufsteigenden Knoten wird ein solcher Durchschnitt in 925 Jahren erfolgen, und
so wird, wie jetzt die Pallas-Bahn in beiden Knoten innerhalb der Ceres-Bahn liegt, nach 1000 Jahren
die Ceres-Bahn innerhalb der Pallas-Bahn liegen. Doch können diese Betrachtungen zu nichts entscheidendem führen, bis die Perturbationen aller drei Bahnen völlig entwickelt sein werden.“

\Needspace{24\baselineskip}
Dr. \textsc{Gauss} fährt indessen, bis dieses geschehen kann, fort, die elliptischen Elemente der Bahn
zu verbessern, und sie mit den besten Meridian-Beobachtungen zu vergleichen. Hier folgt eine solche
Vergleichung dieser sämmtlichen Beobachtungen mit seinen zweiten Elementen der \emph{Juno}.

\begin{center}
\begingroup\fontsize{8}{10}\selectfont\setlength{\tabcolsep}{3pt}\renewcommand{\arraystretch}{1.1}
\begin{tabular}{r|r|r|r|r|l}
 & \multicolumn{2}{c|}{Berechnete} & \multicolumn{2}{c|}{Differenz} & \\ 1804 & \shortstack{Gerade Aufstei-\\gung} & \shortstack{südliche Ab-\\weichung} & \shortstack{in der\\AR.} & \shortstack{in der\\Decl.} & Beobachter \\ \hline
Sept. 13 & $0^\circ\ 44'\ 56''0$ & $1^\circ\ 52'\ 37''6$ & $-0''4$ & $+0''6$ & von Zach \\
14 & $0\ 35\ 42.9$ & $2\ 5\ 40.3$ & $+1.9$ & $+4.8$ & — \\
15 & $0\ 26\ 19.8$ & $2\ 18\ 48.6$ & $-0.2$ & $+10.1$ & — \\
17 & $0\ 7\ 8.0$ & $2\ 45\ 18.8$ & $-0.1$ & $-5.7$ & — \\
18 & $359\ 57\ 20.9$ & $2\ 58\ 38.7$ & $+0.4$ & $+7.8$ & — \\
20 & $359\ 37\ 39.0$ & $3\ 25\ 24.7$ & $.\ .\ .\ .$ & $+21.2$ & — \\
23 & $359\ 7\ 17.4$ & $4\ 5\ 36.5$ & $+2.5$ & $+2.0$ & — \\
23 & $359\ 7\ 3.3$ & $4\ 5\ 55.0$ & $+3.3$ & $+8.0$ & Dr. Burckhardt \\
25 & $358\ 46\ 45.8$ & $4\ 32\ 41.4$ & $-2.2$ & $+13.2$ & Dr. Maskelyne \\
27 & $358\ 26\ 56.9$ & $4\ 58\ 43.7$ & $-3.7$ & $-12.8$ & von Zach \\
28 & $358\ 16\ 58.7$ & $5\ 11\ 49.3$ & $-7.9$ & $+5.5$ & — \\
29 & $358\ 6\ 48.0$ & $5\ 25\ 11.7$ & $-4.3$ & $+15.0$ & Dr. Maskelyne \\
30 & $357\ 57\ 18.2$ & $5\ 37\ 50.4$ & $-8.5$ & $+6.7$ & von Zach \\
Oct. 2 & $357\ 38\ 6.4$ & $6\ 2\ 59.7$ & $-9.6$ & $+0.2$ & — \\
4 & $357\ 19\ 34.8$ & $6\ 27\ 40.2$ & $-4.7$ & $+1.5$ & — \\
5 & $357\ 10\ 37.3$ & $6\ 39\ 43.7$ & $-4.2$ & $+3.8$ & — \\
6 & $357\ 1\ 52.8$ & $6\ 51\ 35.6$ & $-3.5$ & $+3.3$ & — \\
\end{tabular}\endgroup
\end{center}

\SourceLeaf{263}
\indent Nach denselben Elementen hat Dr. \textsc{Gauss} ferner für den künftigen Lauf der \emph{Juno} für Seeberger
Mitternacht folgende Ephemeride berechnet: [vom Oct. 15 bis Dec. 2 enthält sie ger. Aufsteigung und
Abweichung.]

Ohngefähr den 18. December kommen Juno und Ceres geocentrisch nahe zusammen.

\noindent—\quad—\quad—\par

\hfill[\textsc{von Zach.}]

\Needspace{10\baselineskip}
\begin{center}
BEOBACHTETE STERNBEDECKUNG.
\end{center}

\begin{center}
\rule{0.6\linewidth}{0.4pt}

Monatliche Correspondenz. Band X. S. 481. 1804 November.

\rule{0.6\linewidth}{0.4pt}
\end{center}

In Braunschweig beobachtete Dr. \textsc{Gauss} den 16. October 1804, die Bedeckung von $\lambda$ in den Fischen:

\begin{center}
\begingroup\fontsize{10}{12}\selectfont\setlength{\tabcolsep}{3pt}\renewcommand{\arraystretch}{1.1}
\begin{tabular}{llrll}
den Eintritt & um & $10^{\mathrm u}\ 31^{\mathrm m}\ 16^{\mathrm s}7$ & mittl. & Zeit \\
den Austritt & „ & $11\ 40\ 48.0$ & „ & „ \\
\end{tabular}\endgroup
\end{center}
